%--------------------------
\begin{frame}{Question 3}

\begin{block}{}
L’office récepteur est seul compétent pour décider :

\begin{itemize}
  \item[$\bullet$] a. d’exiger une taxe de transmission pour les demandes PCT déposées auprès de lui. 
  \item[$\bullet$] b. d’accepter une seule langue ou plusieurs langues dans laquelle ou dans lesquelles 
  des demandes PCT peuvent être déposées auprès de lui. 
  \item[$\bullet$] c. du choix des administrations chargées de la recherche internationale supplémentaire 
  auxquelles les déposants auront accès s’ils déposent auprès de lui. 
  \item[$\circ$] d. Aucune des réponses a, b ou c n’est correcte. 
\end{itemize}

\end{block}

\end{frame}

%--------------------------
\begin{frame}{Question 3}

\textbf{Base juridique :} \pctrule{14}.1(c) + \pctrule{12}.1(a) + \pctrule{45bis}.9.  

\begin{itemize}
  \item[$\bullet$] \textbf{a.} \; Correct : chaque office récepteur décide d’exiger ou non
  une taxe de transmission (\pctrule{14}.1(c)).

  \item[$\bullet$] \textbf{b.} \; Correct : la langue de dépôt est déterminée par ce que
  l’office récepteur accepte (\pctrule{12}.1(a)).

  \item[$\bullet$] \textbf{c.} \; Correct : l’accès aux administrations chargées de la recherche
  supplémentaire dépend de l’office récepteur compétent (\pctrule{45bis}.9).

  \item[$\circ$] \textbf{d.} \; Faux puisque \textbf{a}, \textbf{b} et \textbf{c} sont correctes.
\end{itemize}

\end{frame}
